
\documentclass[../../main.tex]{subfiles}
\begin{document}

% 年度の大きな表示
\begin{center}
  {\fontsize{48pt}{58pt}\selectfont \textbf{1995}\normalsize\textbf{年度}}
\end{center}

\vspace{10mm}

% 本文


\vspace{15mm}

% 出題テーマと難易度の表
\noindent\textbf{出題テーマと難易度}

\vspace{5mm}

\begin{center}
  \shadowbox{%
    \begin{tabular}{|c|p{8cm}|c|}
      \hline
      \textbf{問題番号} & \textbf{テーマ} & \textbf{難易度} \\
      \hline
      第1問 & 数列の極限と整数条件（階乗を含む数列） & \\
      \hline
      第2問 & 五角形の面積の最大値（外角と三角関数） & \\
      \hline
      第3問 & 関数の増減とグラフの概形，極限の証明 & \\
      \hline
      第4問 & 確率（期待値の計算） & \\
      \hline
    \end{tabular}%
  }
\end{center}

\vspace{15mm}

% 解答
\noindent\textbf{解答}

\vspace{5mm}

\begin{center}
  \shadowbox{%
    \renewcommand{\arraystretch}{1.6}
    \begin{tabular}{|c|c|p{8cm}|}
      \hline
      \textbf{問題番号} & \textbf{小問} & \textbf{答え} \\
      \hline
      \multirow{3}{*}{第1問} & (1) & $\displaystyle\lim_{n\to\infty}a(n)=0$ \\
      \cline{2-3}
                             & (2) & $n=1,2,3,4,6$ \\
      \cline{2-3}
                             & (3) & $n=1,2,3,4,5,6,7$ \\
      \hline
      第2問 & - & 最大値 $\sqrt{2}+1$（$\theta=\dfrac{\pi}{4}$） \\
      \hline
      \multirow{2}{*}{第3問} & (1) & 増減表・概形図示（略） \\
      \cline{2-3}
                             & (2) & 証明問題 \\
      \hline
      \multirow{2}{*}{第4問} & (1) & $E=\dfrac{(n+1)(3n+2)}{12}$ \\
      \cline{2-3}
                             & (2) & $\dfrac{1}{8}$ \\
      \hline
    \end{tabular}%
    \renewcommand{\arraystretch}{1.0}
  }
\end{center}

\vspace{15mm}

% 解答時間
\noindent\textbf{試験形態}

\vspace{5mm}

\begin{center}
  \shadowbox{%
    \begin{tabular}{p{8cm}}
      （要確認）
    \end{tabular}%
  }
\end{center}

\end{document}
